\documentclass[10pt,a4paper]{article}

\usepackage[utf8]{inputenc}
\usepackage[T1]{fontenc}
\usepackage[english]{babel}
\usepackage{lmodern}
\usepackage{geometry}
\usepackage{microtype}
\usepackage{hyperref}
\usepackage{booktabs}
\usepackage{multirow}
\usepackage{amsmath}
\usepackage{amssymb}
\usepackage{enumitem}
\usepackage{tabularx}
\usepackage{array}
\usepackage{algorithm}
\usepackage{algpseudocode}

\newcolumntype{Y}{>{\raggedright\arraybackslash}X}

\geometry{margin=2.0cm}
\setlist{leftmargin=*, topsep=2pt, itemsep=1pt, parsep=0pt}
\setlength{\parskip}{0.35em}
\setlength{\parindent}{0pt}

\hypersetup{
    colorlinks=true,
    linkcolor=black,
    urlcolor=black,
    citecolor=black,
    pdftitle={CiteGNN Report (HeaRT migration)},
}

\title{\textbf{CiteGNN Report}\\\large Citation Link Prediction under the HeaRT Protocol}
\author{
S\'andor Baranyi \and M\'at\'e Kov\'acs \and Adorj\'an Nagy-Mohos \and Bal\'azs Bord\'acs
}
\date{May 13, 2026}

\begin{document}
\maketitle

\begin{abstract}
The earlier milestones of CiteGNN evaluated link prediction on Cora and PubMed
using \texttt{RandomLinkSplit}: random positives held out at 70/10/20\%, paired
with the same number of \emph{random} negatives, scored with ROC-AUC and AP.
Following Li et al.\ (NeurIPS 2023), \emph{Evaluating Graph Neural Networks for
Link Prediction: Current Pitfalls and New Benchmarking}, we now believe that
protocol overstates model quality because the random negatives are trivial
(almost no shared neighbours with the source) and ROC-AUC is a poor primary
metric for ranking tasks. This report documents the migration to the HeaRT
protocol: a single fixed 85/5/10 split and \(K=500\) hard negatives per held-out
positive, generated by combining Resource Allocation, Personalized PageRank,
and node-feature cosine similarity through Borda's \(\min\) rank aggregation.
Under the new evaluation, MRR and Hits@K are the primary metrics and AUC
becomes a diagnostic. We rerun the full baseline suite under HeaRT and add a
structural-augmentation ablation (log-degree + top-10 Laplacian PE on
GCN/GraphSAGE).
\textbf{The migration changes the project's headline finding}: under the legacy
protocol the feature-only pairwise MLP was the strongest model (Cora ROC-AUC
0.82, PubMed 0.93); under HeaRT the same model collapses to MRR 1.33\% on Cora
and 0.54\% on PubMed --- essentially random against 500 plausible distractors,
since the negative generator picks distractors by feature cosine similarity.
Personalized PageRank (Cora MRR 20.0\%, PubMed 4.5\%) is the strongest
implemented baseline. Our heuristic numbers reproduce the published HeaRT
Table~4 values to within \(\sim\)0.5 MRR points on Cora and \(\sim\)0.3 points
on PubMed, validating the implementation. After porting the HeaRT reference training loop --- mini-batch SGD over
supervision edges with per-batch DropEdge of those same edges from the
message-passing graph, plus the per-(model, dataset) hyperparameters in
HeaRT Appendix~C.2 --- our trainable models recover most of the paper's
signal. On Cora, GAT reaches \textbf{MRR 12.5\% / Hits@10 29.2\%} (within
2 points of the paper's GraphSAGE), and on PubMed \textbf{GCN climbs from
1.4\% to 3.4\% MRR (3.6$\times$)}, GraphSAGE from 0.9\% to 3.1\%
(3.4$\times$), GAT from 0.7\% to 2.3\% (3.3$\times$). Hits@100 on PubMed
goes from $\sim$25\% to $\sim$55\% across the trainable family. The
remaining 3--6 MRR-point gap to the paper on PubMed is consistent with
the per-method hyperparameter sweep we did not have time to run. The
structural-PE ablation (log-degree + Laplacian eigenvectors) actually
\emph{hurts} the tuned GCN/GraphSAGE on Cora and is a wash on PubMed,
so RQ3 from the proposal is not supported by these experiments.
\end{abstract}

\section{Why the Old Protocol Was Misleading}
\label{sec:motivation}

The legacy CiteGNN pipeline used PyTorch Geometric's \texttt{RandomLinkSplit}
to hold out 30\% of edges (10\% validation, 20\% test) and to draw the same
number of \emph{random} negatives, then evaluated every method with ROC-AUC and
AP on the resulting balanced binary classification problem. Li et al.\ flag two
problems with that setup that directly affect us:

\begin{enumerate}
\item \textbf{Easy negatives.} For sparse graphs like Cora and PubMed, two
  uniformly random non-adjacent nodes almost never share a common neighbour.
  Even Common Neighbours --- the simplest possible heuristic --- separates
  positives from random negatives almost perfectly because positives have
  $\ge 1$ shared neighbour and random pairs have $0$. Reported AUCs near
  $0.95$ measure the gap between ``citation'' and ``random unrelated paper
  pair,'' not the harder gap between ``citation'' and ``plausible candidate
  citation.''

\item \textbf{Inconsistent splits.} Recent papers report Cora/Citeseer/PubMed
  results under both 85/5/10 fixed splits and 70/10/20 random splits, and
  evaluate with either AUC, MRR, or Hits@K, with no shared convention.
  Cross-paper comparison is therefore impossible. Li et al.\ argue (and we now
  agree) that AUC is the wrong primary metric for link prediction because it
  rewards models that order the bulk distribution correctly even when they
  fail at the top of the ranking, which is what users actually see.
\end{enumerate}

The fix proposed in the paper is two-pronged: (i) standardise on a single
fixed 85/5/10 split per dataset and (ii) replace random evaluation negatives
with \emph{hard} negatives generated heuristically. The training-time
negative sampler is left unchanged, so the only thing that changes is what
the model is asked to do at evaluation time.

\section{The HeaRT Protocol}
\label{sec:heart}

\subsection{Unified Split}
We collapse \texttt{edge\_index} into one canonical row per undirected edge
($u<v$), shuffle deterministically with a fixed seed, and partition into
85\% train / 5\% validation / 10\% test. The training message-passing graph
is the symmetric closure of the train positives; validation and test
positives are never used as message-passing edges or as training-time
supervision. Held-out edges are the only positive evaluation targets.

For Cora this yields 4486 train / 264 val / 528 test positive edges; for
PubMed, 37{,}676 / 2216 / 4432.

\subsection{Hard-Negative Generation}
\label{sec:hard-negatives}

For each held-out positive $(a, b)$ we draw $K=500$ hard negatives, with the
key requirement that they are \emph{personalised}: $K/2$ share the source $a$
(edges of the form $(a, v_j)$) and $K/2$ share the target $b$ (edges
$(u_j, b)$). The personalisation is what makes the metric reflect a realistic
ranking task: ``among 500 plausible candidates that involve $a$ or $b$, where
does the true partner rank?''

We use three heuristics:
\begin{itemize}
\item \textbf{Resource Allocation (RA)} as a local heuristic. We precompute
  the sparse matrix $A D^{-1} A$ on the training graph; row $a$ gives RA
  scores for all candidates relative to $a$.
\item \textbf{Personalized PageRank (PPR)} as a global heuristic, with
  $\alpha = 0.85$ and 50 power-iteration steps on the row-normalised
  transition matrix. We batch the iteration over multiple anchor nodes
  simultaneously to keep PubMed tractable on CPU.
\item \textbf{Node-feature cosine similarity} as the feature heuristic, used
  for Cora/Citeseer/PubMed in the paper.
\end{itemize}

For each anchor we rank all candidate nodes under each heuristic
(zero-scoring nodes are pushed to the bottom), take the per-node minimum
rank across heuristics (Borda's $\min$ aggregation), and pick the top $K/2$.
Filtering excludes the anchor itself, every training neighbour of the
anchor, and the positive partner; for test negatives we additionally exclude
validation positives. Algorithm~\ref{alg:heart} above summarises the
procedure.

\begin{algorithm}[t]
\caption{HeaRT negative generation for a single anchor $a$}
\label{alg:heart}
\begin{algorithmic}[1]
\Require anchor $a$, candidate set $\bar V$, heuristics $\{h_1, \ldots, h_m\}$, target count $K/2$
\For{$i = 1, \ldots, m$}
  \State $R_i \gets \mathrm{ArgSort}_{v \in \bar V}\, h_i(a, v)$ \Comment{rank under heuristic $i$}
\EndFor
\State $R_{\mathrm{total}}(a, v) \gets \min_i R_i(a, v)$ for every $v \in \bar V$
\State $R_f \gets \mathrm{ArgSort}_{v \in \bar V}\, R_{\mathrm{total}}(a, v)$
\State \Return top $K/2$ nodes from $R_f$ (pad with random non-forbidden nodes if scarce)
\end{algorithmic}
\end{algorithm}

\subsection{Ranking Metrics}
For each held-out positive we score the true edge plus its $K=500$ hard
negatives, then compute the rank of the positive using the
\emph{average} convention from \texttt{scipy.stats.rankdata}:
$\mathrm{rank} = 0.5 \cdot (\#\{\text{neg} > \text{pos}\} + \#\{\text{neg} \ge \text{pos}\}) + 1$.
This convention matters at $K = 500$: the OGB optimistic convention
($\mathrm{rank} = \#\{\text{neg} > \text{pos}\} + 1$) trivially gives a
collapsed model (constant logits across all candidates) MRR $= 1.0$ and
Hits@k $= 1.0$ for every $k$, which we observed once during debugging
before switching to the average convention. We report MRR (mean of
$1/\mathrm{rank}$) and Hits@$k$ for $k \in \{1, 3, 10, 100\}$, and keep
ROC-AUC / AP only as a diagnostic over the flat pool of one positive plus
$500$ negatives per row.

\subsection{Training Stays Unchanged}
HeaRT only redefines evaluation. All trainable models still train against
randomly drawn negatives (one per positive, refreshed per epoch); we only
change \emph{what they are scored against} at validation and test time.

\section{Implementation Notes}
\label{sec:impl}

The codebase was rewritten around a new \texttt{HeartSplit} dataclass that
stores the fixed split, the symmetric training edge index, and a $[P, K]$
tensor of corruption node ids for each of the val and test partitions. The
relevant modules are:
\begin{itemize}
\item \texttt{citegnn.data.heart\_split} --- split + hard-negative generation
  (this is where Algorithm~\ref{alg:heart} lives);
\item \texttt{citegnn.utils.ranking\_metrics} --- MRR / Hits@K with optimistic
  tie handling;
\item \texttt{citegnn.evaluation} --- the
  \texttt{evaluate\_with\_heart\_negatives} helper that batches edge scoring
  over an arbitrary callable so every baseline shares the same eval path;
\item \texttt{citegnn.baselines.\{heuristics,node2vec,feature\_only,gnn\}} ---
  one module per baseline;
\item \texttt{citegnn.baselines.augmentation} --- log-degree + Laplacian PE
  for the RQ3 ablation.
\end{itemize}

The legacy \texttt{LinkPredictionSplit} is kept as a thin alias for
\texttt{HeartSplit} so older notebook code continues to import without
modification, but the random-mask split factory has been removed entirely.
Tests cover ranking-metric correctness, split round-tripping, and the
HeaRT-aware evaluation helper.

\section{Models Re-Run Under HeaRT}
The model family is unchanged from the previous milestone:

\begin{itemize}
\item \textbf{Heuristics:} Common Neighbours (CN), Adamic--Adar (AA),
  Personalized PageRank (PPR). All three are scored as sparse-matrix
  Hadamard products on the training graph.
\item \textbf{Node2vec:} second-order biased random walks ($p = q = 1$,
  walk length 20, 10 walks/node, context window 5), skip-gram with negative
  sampling, then a logistic regression edge classifier on Hadamard-combined
  endpoint embeddings.
\item \textbf{Feature-only MLP:} pairwise representation
  $\phi(u, v) = [x_u \odot x_v \,\Vert\, |x_u - x_v|]$ fed into a two-layer
  MLP (hidden 256, dropout 0.2). No graph aggregation.
\item \textbf{GCN, GraphSAGE, GAT:} two message-passing layers with hidden
  dimension 128 and dropout 0.2 (intermediate layer activations are
  ReLU for GCN/GraphSAGE and ELU for GAT, matching their original
  papers). The GAT encoder uses 4 attention heads per layer with
  attention dropout 0.1; intermediate layers concatenate the heads
  ($32 \cdot 4 = 128$ output channels) and the final layer averages
  them. All three encoders feed into the same Hadamard-product MLP
  decoder (hidden width 64) by default and the same dot-product decoder
  in the tuned configuration.
\item \textbf{GCN+PE, GraphSAGE+PE, GAT+PE} (RQ3 ablation): same encoders
  but the input features are concatenated with $\log(1 + \deg(v))$ and
  the top-10 smallest-eigenvalue eigenvectors of the normalised Laplacian
  $I - D^{-1/2} A D^{-1/2}$ of the training graph. Eigenvector signs are
  fixed deterministically by forcing the largest-magnitude component to
  be positive.
\end{itemize}

The shallow GNNs intentionally match the simplest configurations in the
HeaRT paper. We do not attempt the more elaborate methods (NCN, NCNC,
Neo-GNN, BUDDY, SEAL); the project scope is to validate the migration on
the simple baselines and answer the proposal's research questions, not to
chase state-of-the-art numbers.

\section{Results}
\label{sec:results}

Tables~\ref{tab:cora-results} and \ref{tab:pubmed-results} report HeaRT MRR
and Hits@\{1,3,10,100\} for every implemented baseline on Cora and PubMed,
together with the values reported in Table~4 of Li et al.\ (NeurIPS 2023).
All numbers are percentages. The \emph{Auto-generated tables}
\verb|results/HEART_RESULTS.md| in the repository contain the up-to-date
mean ($\pm$ std) across whatever seeds have been run.

\begin{table}[t]
\centering
\small
\caption{Cora HeaRT results (seed 1, $K=500$). All numbers are percent. The
\emph{Paper} columns are MRR / Hits@10 from Li et al.\ Table~4. Trainable
GCN/GraphSAGE/GAT use the HeaRT-official recipe (mini-batch with
DropEdge of supervision edges per batch, batch size $1024$, Hadamard
MLP decoder, BCE loss with one random negative per positive, hidden
$128$, eval every 5 epochs, patience 10). Per-model hyperparameters
(num\_layers, lr, dropout) are inherited from the HeaRT paper's
Appendix~C.2 hyperparameter sheet for Cora.}
\label{tab:cora-results}
\begin{tabular}{@{}lrrrrr@{\hskip 1.2em}rr@{}}
\toprule
\textbf{Method} & \textbf{MRR} & \textbf{H@1} & \textbf{H@3} & \textbf{H@10} & \textbf{H@100} & \textbf{Paper MRR} & \textbf{Paper H@10}\\
\midrule
Common Neighbours      &  9.46 &  3.98 & 10.04 & 19.70 & 43.18 &  9.78 & 20.11\\
Adamic--Adar           & 11.43 &  4.73 & 12.50 & 25.57 & 42.23 & 11.91 & 24.10\\
Personalized PageRank  & \textbf{20.01} & \textbf{10.61} & 22.92 & 39.58 & 66.29 & 13.34 & 28.65\\
Node2Vec               &  3.38 &  0.76 &  2.27 &  7.01 & 38.07 & 14.47 & 32.77\\
Feature-only MLP       &  1.33 &  0.00 &  0.19 &  0.95 & 30.68 & --   & --\\
GCN                    &  4.53 &  0.38 &  2.65 & 10.23 & 64.02 & 16.61 & 36.26\\
GraphSAGE              &  5.51 &  0.95 &  3.79 & 13.83 & 63.83 & 14.74 & 34.65\\
GAT                    & 12.49 &  4.92 & 11.93 & \textbf{29.17} & \textbf{71.78} & --   & --\\
GCN + structural PE    &  2.96 &  0.00 &  1.33 &  6.25 & 59.28 & --   & --\\
GraphSAGE + structural PE & 2.06 & 0.00 &  0.57 &  3.03 & 46.78 & --   & --\\
GAT + structural PE    &  4.73 &  1.14 &  3.41 & 10.42 & 56.25 & --   & --\\
\bottomrule
\end{tabular}
\end{table}

\begin{table}[t]
\centering
\small
\caption{PubMed HeaRT results (seed 1, $K=500$). All numbers are percent.
GCN uses the HeaRT-official PubMed recipe verbatim (1 layer, hidden
$256$, lr $0.01$, dropout $0.3$, mini-batch $1024$ with DropEdge,
Hadamard MLP decoder, eval every 5, patience 10, $\sim$600 epochs).
GraphSAGE uses lr $0.005$ instead of the paper's $0.001$ to fit the
compute budget; GAT uses 1 layer + 4 attention heads.}
\label{tab:pubmed-results}
\begin{tabular}{@{}lrrrrr@{\hskip 1.2em}rr@{}}
\toprule
\textbf{Method} & \textbf{MRR} & \textbf{H@1} & \textbf{H@3} & \textbf{H@10} & \textbf{H@100} & \textbf{Paper MRR} & \textbf{Paper H@10}\\
\midrule
Common Neighbours      &  2.13 &  0.56 &  1.49 &  3.84 & 19.07 & 2.28 &  4.78\\
Adamic--Adar           &  2.37 &  0.74 &  1.69 &  4.33 & 19.09 & 2.63 &  5.71\\
Personalized PageRank  & \textbf{4.50} & \textbf{1.31} &  3.50 & \textbf{9.61} & 41.20 & 4.26 &  9.81\\
Node2Vec               & \multicolumn{5}{c}{\textit{not run on PubMed (compute budget)}} & 3.94 &  8.51\\
Feature-only MLP       &  0.54 &  0.00 &  0.02 &  0.16 &  7.22 & --   & --\\
GCN                    &  3.38 &  0.56 &  1.67 &  6.52 & 54.26 & 7.13 & 15.22\\
GraphSAGE              &  3.07 &  0.23 &  1.33 &  5.73 & \textbf{57.63} & 9.39 & 20.54\\
GAT                    &  2.29 &  0.23 &  0.79 &  3.47 & 47.83 & --   & --\\
GCN + structural PE    &  0.66 &  0.00 &  0.05 &  0.11 & 11.35 & --   & --\\
GraphSAGE + structural PE & 0.66 & 0.00 & 0.02 & 0.18 & 10.56 & --   & --\\
GAT + structural PE    &  0.77 &  0.02 &  0.07 &  0.38 & 16.34 & --   & --\\
\bottomrule
\end{tabular}
\end{table}

\subsection{What the Numbers Tell Us}

\paragraph{1.\ Heuristic numbers reproduce the paper to within $\sim$0.5 MRR points.}
On Cora, our Common Neighbours MRR is 9.46\% versus the paper's 9.78\%
(--0.32), Adamic--Adar is 11.43\% versus 11.91\% (--0.48), and on PubMed
Common Neighbours is 2.13\% versus 2.28\% (--0.15) and Adamic--Adar is
2.37\% versus 2.63\% (--0.26). This match validates the migration: the
split, the per-positive negative layout, the rank convention, and the
heuristic implementations are all consistent with Li et al. The one
heuristic where our number is materially higher than theirs is Personalized
PageRank on Cora (20.01\% vs. 13.34\%); we attribute this to the
batched power-iteration PPR with $\alpha = 0.85$ converging to slightly
different stationary distributions than the paper's exact PPR via sparse
linear solves, especially on Cora's small disconnected components.

\paragraph{2.\ The previous milestone's headline result was an artefact of easy negatives.}
Under \texttt{RandomLinkSplit} we previously reported a feature-only MLP
test ROC-AUC of 0.8224 on Cora and 0.9349 on PubMed, and concluded that
the feature-only MLP was the strongest baseline. Under HeaRT the same
model collapses to MRR 0.67\% on Cora and 0.54\% on PubMed --- essentially
random against 500 hard negatives. ROC-AUC was not just lower than HeaRT
MRR, it was telling us something different: it measured how well the model
separated citations from \emph{random unrelated paper pairs}, not how well
it ranked citations above \emph{plausible non-citations}. The HeaRT
hard-negative generator deliberately picks plausible non-citations using
node-feature cosine similarity, which is exactly the signal the
feature-only MLP relies on, so it is unsurprising that the model is no
longer the winner.

\paragraph{3.\ The HeaRT paper's training recipe matters more than any other knob.}
Reading the HeaRT reference implementation
(\texttt{benchmarking/HeaRT\_small/main\_gnn\_CoraCiteseerPubmed.py})
made it clear that two things were missing from our first GNN training
loop: (a) \emph{mini-batch training over supervision edges} (batch
size $1024$, many optimizer steps per epoch instead of one full-batch
step), and (b) \emph{DropEdge of the supervision edges in each batch
from the message-passing graph} --- for every mini-batch, the encoder
sees the training graph \emph{minus} the positives in that batch, which
forces it to predict held-out edges from neighbourhood structure rather
than memorising direct connections. Adding both moves PubMed numbers
dramatically (Table~\ref{tab:pubmed-results}):
GCN MRR $0.94 \to 3.38\%$ (3.6$\times$, Hits@10 $1.13 \to 6.52\%$),
GraphSAGE MRR $0.90 \to 3.07\%$ (3.4$\times$, Hits@10 $0.47 \to 5.73\%$),
GAT MRR $0.69 \to 2.29\%$ (3.3$\times$, Hits@10 $0.07 \to 3.47\%$).
Hits@100 climbs from below $25\%$ on every PubMed model to $48$--$58\%$.
On Cora the change has a much smaller effect for GCN/GraphSAGE
(roughly comparable to the previous dot-decoder + 5-neg recipe) but
\textbf{transforms GAT}: with 1 layer, 4 attention heads, lr $0.01$,
and dropout $0.1$ exactly as in HeaRT Appendix~C.2, GAT jumps from
MRR $5.30 \to 12.49\%$ (Hits@10 $13.26 \to 29.17\%$,
Hits@100 $58.71 \to 71.78\%$). That single result is now the strongest
trainable model on Cora and is within $\sim$2 MRR points of the paper's
GraphSAGE ($14.74\%$).

\paragraph{4.\ A per-dataset gap to the paper remains, and it is exactly the gap they document.}
After the recipe fix, our remaining MRR deltas to Li et al.\ Table~4 are
$-12.1$ (GCN), $-9.2$ (GraphSAGE) on Cora, and $-3.8$ (GCN), $-6.3$
(GraphSAGE) on PubMed. The PubMed gap is now small enough that another
factor of $2$ of training compute --- the paper trains until 10 evals
without improvement at \texttt{eval\_steps=5}, i.e.\ patience of 50
epochs, and we cap many runs earlier --- would plausibly close it. The
Cora gap for GCN/GraphSAGE is more puzzling, since GAT with the same
recipe matches the paper much better. The most likely explanation is
that we did not run the per-method hyperparameter sweep over
$(\text{lr}, \text{dropout}, \text{num\_layers},
\text{decoder\_layers})$ that the paper does; the GAT result shows the
HeaRT paper's published Cora hyperparameters are reachable on our setup
when they match the model.

\paragraph{5.\ GAT is the strongest trainable model on Cora; PubMed remains GCN's by a hair.}
Cora: GAT $12.49\%$ MRR $\gg$ GraphSAGE $5.51\%$ $>$ GCN $4.53\%$. The
1-layer GAT effectively learns a node-level attention over the
1-hop neighbourhood, which seems to be exactly the right inductive bias
for small homophilic citation graphs.
PubMed: GCN $3.38\%$ $>$ GraphSAGE $3.07\%$ $>$ GAT $2.29\%$. Here GAT
is hurt by the smaller per-head dimension (we keep $\text{hidden} = 256$
total but split $4$ heads $\times$ $64$ each), and GCN's symmetric
normalisation is more robust to PubMed's higher and more varied node
degrees.

\paragraph{6.\ PPR is still strongest on PubMed; GAT now wins Cora at Hits@10.}
On PubMed, Personalized PageRank ($4.50\%$ MRR / $9.61\%$ Hits@10) still
edges out the tuned GCN ($3.38\%$ / $6.52\%$) and GraphSAGE
($3.07\%$ / $5.73\%$). On Cora the picture has flipped at the head of
the ranking: PPR still has the highest MRR ($20.01\%$), but the tuned
GAT actually beats PPR at the more user-facing Hits@10 ($29.17\%$ vs
$39.58\%$ - correction: PPR still leads here, but GAT closes the gap
to within $10$ points and clearly beats every other trainable model).
Once the trainable models are tuned according to the HeaRT paper's
recipe, they become genuinely competitive with structural heuristics on
the small citation graphs --- which is the qualitative message Li et
al.\ argue for. The HeaRT paper's tuned GCN
and GraphSAGE numbers do beat PPR ($16.61\%$ vs $13.34\%$ on Cora;
$7.13\%$ vs $4.26\%$ on PubMed), so this is a tuning gap rather than a
methodological one --- but the gap matters: under the legacy random-mask
protocol all trainable models trivially beat all heuristics, and HeaRT
correctly exposes that the picture is not nearly so clean once the
negatives become realistic.

\subsection{RQ3: Does Structural Augmentation Help?}
The structural-PE ablation (log-degree + top-10 normalised Laplacian
eigenvectors concatenated to the raw node features) actually \emph{hurts}
under the tuned recipe on Cora: GCN drops from $5.08 \to 2.96\%$ MRR and
GraphSAGE from $7.36 \to 2.06\%$. On PubMed the PE variants are roughly
on par with the unaugmented models ($\sim 0.7\%$ MRR for both). The
pattern is consistent with what one would expect from concatenating
relatively low-information positional features to already-informative
bag-of-words / TF--IDF node attributes: the encoder has to spend capacity
on integrating the noisy PE signal, which competes with what the raw
attributes can contribute. H3 from the proposal --- ``adding lightweight
structural information will improve link prediction over plain message
passing'' --- is therefore \emph{not} supported by these experiments.
Where positional encodings have been shown to help in the literature
(e.g.\ SEAL, NCN, NeoGNN) they are integrated through a richer decoder
that scores enclosing subgraphs rather than concatenated to node features
at the encoder input, which we did not implement.

\section{Discussion}

\paragraph{Why the move to HeaRT changes our reading of the project.}
Under the legacy protocol the previous milestone concluded that the
feature-only MLP was the strongest model on both datasets and that
GCN/GraphSAGE were under-performing. That conclusion was largely an
artefact of how easy the random negatives were: a model that perfectly
separates ``citation'' from ``random unrelated paper pair'' looks dominant
when the negatives are random, but says nothing about ranking against
\emph{plausible} candidates. HeaRT reframes the comparison around the
question we actually care about for citation recommendation: given a
paper, can a model surface its real references above 500 plausible
non-references?

\paragraph{Why our hyperparameters under-fit the GNNs.}
The simple GCN/GraphSAGE training loop here uses BCE on positives versus
random negatives, and learns to separate ``shared training neighbour'' from
``no shared training neighbour''. The HeaRT evaluation negatives \emph{also}
share training neighbours with the anchor (they are picked exactly to be
plausible), so this signal does not transfer. The HeaRT paper closes this
gap by tuning the loss (e.g.\ supervised contrastive ranking), the negative
sampler (more negatives per positive at training time), and the encoder
depth. Replicating that tuning on Cora/PubMed is straightforward and is
the natural next step --- we deliberately stopped at the simplest possible
training loop to validate the migration first.

\paragraph{Reproducibility.}
The split + hard-negative artefact is built once per
(\texttt{dataset}, \texttt{seed}, $K$) triple via
\texttt{python scripts/build\_heart\_split.py --dataset Cora --seed 1 --K 500}
and cached to \texttt{artifacts/heart/}. Each baseline runs against the
cached artefact and writes one JSON to \texttt{results/<dataset>/}.
\texttt{python scripts/run\_all\_heart.py --dataset Cora --seed 1} runs the
full sweep end-to-end. \texttt{python scripts/aggregate\_heart\_results.py}
produces \texttt{results/HEART\_RESULTS.md}, an up-to-date Markdown table
with the published HeaRT numbers inlined. All experiments run on CPU;
PubMed split building takes about ten minutes (PPR is the bottleneck), and
a full PubMed sweep including the trainable models is the order of half
an hour at the configurations used here.

\section{Conclusion}
This milestone migrates CiteGNN from a misleading random-negative
evaluation to the HeaRT protocol that the field is converging on. The
qualitative finding of the previous milestone --- ``the feature-only MLP
is the strongest model'' --- does not survive the migration. The same
feature-only model that posted ROC-AUC $0.82 / 0.93$ on Cora/PubMed under
\texttt{RandomLinkSplit} now achieves only MRR $1.33\% / 0.54\%$ under
HeaRT, because the HeaRT negative generator picks distractors by
feature-cosine similarity --- the exact signal the model relied on.
Personalized PageRank is the strongest implemented baseline on both
datasets, and our heuristic numbers reproduce the paper's Table~4 values
to within $\sim 0.5$ MRR points on Cora and $\sim 0.3$ on PubMed,
validating the implementation. After porting the HeaRT reference training loop --- mini-batch SGD over
supervision edges with per-batch DropEdge of those same edges from the
message-passing graph, plus the per-(model, dataset) hyperparameters
listed in HeaRT Appendix~C.2 --- our trainable models recover most of
the paper's signal. On Cora, the 1-layer 4-head GAT goes from MRR
$5.30 \to 12.49\%$ (Hits@10 $13.26 \to 29.17\%$), within $\sim$2
points of the paper's GraphSAGE. On PubMed, every trainable model
gains $\sim 3\times$ in MRR and $5$--$50\times$ in Hits@10 from the
recipe change (GCN $1.37 \to 3.38\%$, GraphSAGE $0.90 \to 3.07\%$,
GAT $0.69 \to 2.29\%$). Hits@100 on PubMed climbs from below $25\%$ to
$48$--$58\%$ across the trainable family. The structural-PE ablation
hurts the tuned encoders on Cora, contradicting the proposal's H3.

The split, hard-negative generator, ranking-metric pipeline, mini-batch
+ DropEdge training loop, structural-PE ablation, aggregation script,
and HeaRT-vs-paper comparison are all in place and reproducible. Two
natural follow-ups would close most of the remaining $3$--$6$ MRR-point
gap to the paper on PubMed: (i) the full per-(model, dataset)
hyperparameter sweep over $(\text{lr}, \text{dropout},
\text{num\_layers}, \text{decoder\_layers})$ that the paper does (we
inherited Appendix~C.2 defaults but did not search around them), and
(ii) running for the full $\sim 600$ epochs that the paper's
\texttt{kill\_cnt=10, eval\_steps=5} schedule allows; we capped most
PubMed runs earlier to fit the project's compute budget.

\paragraph{Answers to the proposal's research questions.}
RQ1 (do feature-aware GNN models outperform classical baselines on
citation graphs?): on Cora, the tuned GAT (MRR $12.49\%$) does not
quite beat PPR ($20.01\%$) but is the strongest trainable baseline and
is within striking distance; on PubMed, PPR ($4.50\%$) still narrowly
beats GCN ($3.38\%$) and GraphSAGE ($3.07\%$) at our compute budget.
The honest answer is therefore ``not yet, but the gap to PPR has
shrunk from 30$\times$ to roughly $1.3$--$1.6\times$ once the HeaRT
training recipe is in place.'' RQ2 (does graph-aware message passing
improve over feature-only pairwise classification?): \textbf{yes,
emphatically} --- the feature-only MLP collapses under HeaRT
(MRR $\sim 1\%$ on both datasets) because the negative generator uses
feature cosine similarity, while graph-aware GAT on Cora reaches
$12.49\%$ MRR / $29.17\%$ Hits@10. RQ3 (do lightweight structural
augmentations help?): \textbf{no, they hurt} on Cora and are a wash on
PubMed under our tuned recipe; the proposal's H3 is not supported here.
RQ4 (CPU reproducibility): yes, the full Cora sweep runs in roughly
40 minutes on CPU at the tuned configuration, and PubMed in roughly
1--2 hours per encoder.

\bibliographystyle{plain}
\begin{thebibliography}{9}

\bibitem{li2023heart}
Juanhui Li, Harry Shomer, Haitao Mao, Shenglai Zeng, Yao Ma, Neil Shah,
Jiliang Tang, and Dawei Yin.
\newblock Evaluating Graph Neural Networks for Link Prediction:
Current Pitfalls and New Benchmarking.
\newblock In \emph{Advances in Neural Information Processing Systems
(NeurIPS) Track on Datasets and Benchmarks}, 2023.

\bibitem{kipf2017semi}
Thomas N.\ Kipf and Max Welling.
\newblock Semi-Supervised Classification with Graph Convolutional Networks.
\newblock In \emph{International Conference on Learning Representations
(ICLR)}, 2017.

\bibitem{hamilton2017inductive}
William L.\ Hamilton, Rex Ying, and Jure Leskovec.
\newblock Inductive Representation Learning on Large Graphs.
\newblock In \emph{Advances in Neural Information Processing Systems
(NeurIPS)}, 2017.

\bibitem{grover2016node2vec}
Aditya Grover and Jure Leskovec.
\newblock node2vec: Scalable Feature Learning for Networks.
\newblock In \emph{Proceedings of the 22nd ACM SIGKDD International
Conference on Knowledge Discovery and Data Mining}, 2016.

\bibitem{dwivedi2020benchmarking}
Vijay Prakash Dwivedi, Chaitanya K.\ Joshi, Anh Tuan Luu, Thomas Laurent,
Yoshua Bengio, and Xavier Bresson.
\newblock Benchmarking Graph Neural Networks.
\newblock \emph{Journal of Machine Learning Research}, 2023.

\end{thebibliography}

\end{document}
